\section{Введение}

\subsection{Актуальность и мотивация}

Графические ускорители NVIDIA GPU занимают центральное место в современных
высокопроизводительных вычислениях: от обучения крупных языковых моделей до
физического моделирования и научных расчётов. Программирование GPU на платформе
CUDA строится на модели SIMT (Single Instruction, Multiple Threads), в которой
группы из 32~потоков~--- варпы~--- исполняют одну и ту же инструкцию над
различными данными~\cite{Lindholm2008, Nickolls2008}. Ключевым промежуточным
представлением в данной экосистеме служит PTX (Parallel Thread Execution)~---
аппаратно-независимый ассемблер NVIDIA~\cite{PTXISA}, в который компилируется
CUDA~C++ и из которого генерируется машинный код для конкретного поколения GPU.

Официальные инструменты NVIDIA (Nsight, Compute Sanitizer) требуют наличия
физического GPU и предоставляют агрегированную статистику на уровне SM,
не обеспечивая детальных метрик на уровне отдельных инструкций PTX.
Программный эмулятор PTX~ISA позволяет устранить оба ограничения: функционировать
на стандартном CPU без GPU-оборудования и собирать произвольные метрики
уровня инструкций.

Настоящий отчёт по практике посвящён двум взаимосвязанным аспектам реализации
такого эмулятора:
\begin{enumerate}
  \item разработке системы анализа производительности для эмулятора CUDA PTX;
  \item разработке тестовых CUDA-ядер и оценке корректности и производительности эмулятора.
\end{enumerate}

\subsection{Цели и задачи}

\textbf{Цель работы}~--- разработать систему профилирования, интегрированную в
программный эмулятор PTX~ISA, и верифицировать корректность эмулятора на
разнообразных реальных CUDA-ядрах.

Для достижения указанной цели поставлены следующие \textbf{задачи}:

\begin{enumerate}
  \item Спроектировать архитектуру системы метрик с кодогенерацией по JSON-манифесту,
        обеспечивающей добавление новой метрики без изменения C++-кода.
  \item Реализовать четыре ключевые метрики производительности:
        эффективность дивергенции ветвлений (\texttt{branch\_efficiency}),
        конфликты банков разделяемой памяти (\texttt{bank\_conflicts}),
        коэффициент слияния транзакций глобальной памяти (\texttt{global\_coalescing})
        и статистику записей в регистровый файл (\texttt{write\_ops}).
  \item Разработать инструмент генерации HTML-отчёта с аннотацией строк
        исходного \texttt{.cu}-файла по данным \texttt{.loc}-директив PTX.
  \item Разработать набор из семи интеграционных тестов на реальных CUDA-ядрах,
        перекрывающих качественно различные вычислительные паттерны.
  \item Провести верификацию корректности эмулятора и оценить производительность
        в сравнении с аппаратным GPU V100-PCIE.
\end{enumerate}

\subsection{Практическая значимость}

Разработанная система профилирования имеет непосредственное практическое применение:

\begin{itemize}
  \item \textbf{Детальное профилирование без GPU.} HTML-отчёт с аннотацией
        строк исходного кода позволяет определить «горячие» инструкции
        и оптимизировать шаблоны обращения к памяти на стандартном CPU.
  \item \textbf{Внедрение в конвейеры CI.} Эмулятор запускается как обычный
        Linux-процесс; для включения в \texttt{GitHub Actions} или
        \texttt{Jenkins} достаточно однострочного скрипта \texttt{cuemu ./binary}.
  \item \textbf{Образовательная ценность.} Метрики дивергенции и конфликтов
        банков обеспечивают наглядную демонстрацию механизмов SIMT-исполнения.
\end{itemize}

\newpage
